\documentclass[UTF8]{ctexart}
\usepackage{geometry}
\usepackage{hyperref}
\usepackage{listings}
\usepackage{xcolor}
\usepackage{booktabs}
\usepackage{graphicx}
\usepackage{amsmath}

\geometry{a4paper, left=2cm, right=2cm, top=2.5cm, bottom=2.5cm}

\hypersetup{
    colorlinks=true,
    linkcolor=blue,
    filecolor=magenta,      
    urlcolor=cyan,
}

\lstset{
    backgroundcolor=\color{lightgray!15},
    basicstyle=\small\ttfamily,
    columns=fullflexible,
    breaklines=true,
    captionpos=b,
    commentstyle=\color{gray},
    keywordstyle=\color{blue},
    frame=single,
    rulecolor=\color{gray!30}
}

\title{\textbf{相识北洋（Meeting in Beiyang）局域网社交系统全景使用与开发手册}}
\author{组长：田华阳 \\ 组员：黄嘉俊 \quad 康猛 \quad 杨浚哲}
\date{\today}

\begin{document}

\maketitle
\tableofcontents
\newpage

\section{引言与系统概述}
“相识北洋（Meeting in Beiyang）”是一款基于 Python 3 与 Flet 0.85.3 框架开发的校园局域网点对点（P2P）社交应用程序。在无中心化服务器的物理局域网（LAN）环境下，该应用利用网络广播与异步套接字通信技术，构建起了一套兼具高并发、自动容错与跨平台能力的局域网社交生态体系。

本手册旨在对系统的技术架构、底层协议、环境安装、编译部署（含 Android 平台容器化打包）、测试模拟及各功能模块的操作细节进行全面深入的拆解与指导，为用户和开发者提供完整的项目指南。

\subsection{课程挑战要求与项目完成情况}
本项目针对课程提出的三项核心挑战进行了完整的设计与实现。以下是各阶段挑战的具体要求与本项目在“相识北洋”系统中的技术实现对照情况：

\subsubsection{课程终极挑战总体要求}
本课程终极挑战要求开发包含 3 个挑战的跨平台局域网社交软件，提供源码、APK 安装包与项目说明书。
本项目在 \texttt{challenge3} 文件夹中对上述要求进行了全功能整合：
\begin{itemize}
    \item \textbf{功能全集成}：在“相识北洋”这一款软件中完美整合了剪贴板分享、大文件断点续传、好友匹配与离线 Gossip 中继等所有核心功能。
    \item \textbf{多平台打包}：编写了防缓存的 Android 自动化编译脚本 \texttt{build-apk-clean.ps1}，用于生成最终的 \texttt{.apk} 安装包。
    \item \textbf{手册撰写}：撰写了本 LaTeX 技术手册并编译成对应的 PDF 版本。
\end{itemize}

\subsubsection{挑战 1：程序分享}
挑战 1 侧重于两端设备间的代码/文本的快速分享与剪贴板集成。
本项目的完成情况如下：
\begin{itemize}
    \item \textbf{剪贴板完美集成}：基于 Flet 的底层机制，在聊天详情页中深度集成了系统剪贴板。消息详情和文件传输卡片中均配备了一键复制文本和文件路径（复制路径）的逻辑。
    \item \textbf{局域网设备自动发现}：设计了 \texttt{UDPService} 服务（监听端口 \texttt{8890}），各端在局域网内定时广播 PING 帧宣告上线，并动态构建局域网节点映射表。
    \item \textbf{跨平台连接}：支持 Windows 桌面端和 Android 移动端的面对面直连通信。
\end{itemize}

\subsubsection{挑战 2：文件速传}
挑战 2 侧重于大文件的传输进度显示、中断重传、压缩包解压与安全哈希校验。
本项目的完成情况如下：
\begin{itemize}
    \item \textbf{自定义保存目录}：集成了 Flet \texttt{FilePicker} 交互，并在设置中允许用户自主配置默认的文件接收目录。
    \item \textbf{ZIP 自动解压}：检测到接收到的文件是 ZIP 压缩包时，消息卡片提供“解压 ZIP”一键后台异步解压缩。
    \item \textbf{分块传输与进度展示}：消息服务将文件切分为 \texttt{256KB} 块进行传输，前端 UI 实时渲染传输进度条、传送速度以及取消动作。
    \item \textbf{断点续传与完整性校验}：传输中断后可基于 part 文件的实际尺寸进行断点续传；传输完成后在 \texttt{FileStore} 中利用 SHA-256 哈希值进行整机安全校验。
\end{itemize}

\subsubsection{挑战 3：相识北洋}
挑战 3 侧重于个人简介条件匹配、好友管理、多跳离线消息中继以及移动状态下对 IP 变更的动态适应。
本项目的完成情况如下：
\begin{itemize}
    \item \textbf{自动审批匹配}：用户可配置兴趣标签与匹配条件，接收申请时自动判定交集标签数，符合条件的自动同意加为好友。
    \item \textbf{好友管理系统}：实现了分组、好友搜索、删除以及实时聊天记录数据库持久化。
    \item \textbf{多跳离线消息中继}：实现了最大 3 跳的 Gossip 中继机制。当接收方离线时，消息自动通过局域网内活跃的在线共同好友（中继节点）暂存并在后续偶遇时投递。
    \item \textbf{动态 IP 漫游支持}：好友关系建立后，若用户切换网络导致 IP 地址改变，可通过后台 UDP 持续广播发现或心跳包（\texttt{HEARTBEAT}）自动更新本地好友的在线 IP 及端口，维持 TCP 会话的连续性。
\end{itemize}

\section{GitHub 仓库获取与安全更新指南}

\subsection{GitHub 仓库地址}
项目托管于 GitHub，官方仓库地址为：
\href{https://github.com/Ortenssia/Meeting-in-Beiyang.git}{\texttt{https://github.com/Ortenssia/Meeting-in-Beiyang.git}}

\subsection{首次克隆与常规更新}
\subsubsection{首次下载}
使用标准 Git 工具克隆仓库到本地：
\begin{lstlisting}[language=bash]
git clone https://github.com/Ortenssia/Meeting-in-Beiyang.git
cd Meeting-in-Beiyang
\end{lstlisting}

\subsubsection{常规代码更新}
当仓库有代码提交时，请在项目根目录下执行以下命令拉取最新代码：
\begin{lstlisting}[language=bash]
git pull origin main
\end{lstlisting}

\subsection{处理代码更新冲突}
如果您在本地对某些代码进行了临时修改，导致 \texttt{git pull} 时提示冲突并失败，请执行以下命令强制覆盖本地代码（这只会覆盖代码，您的本地聊天记录和好友数据库不会丢失）：
\begin{lstlisting}[language=bash]
git fetch --all
git reset --hard origin/main
\end{lstlisting}

\subsection{无 Git 环境下的 ZIP 安全更新步骤}
如果您是直接从 GitHub 网页下载的 ZIP 压缩包，后续更新时请\textbf{千万不要直接用新解压的整个文件夹覆盖原有文件夹}。这会导致您的好友关系和历史消息被彻底抹除。
请按照以下步骤进行安全覆盖：
\begin{enumerate}
    \item 解压最新下载的 ZIP 压缩包；
    \item 将解压出来的 \texttt{core/} 目录以及根目录下的 \texttt{requirements.txt}、\texttt{pyproject.toml} 等代码文件复制并\textbf{覆盖}到您原有的本地目录中；
    \item \textbf{千万不要覆盖或删除}原有本地目录中的 \texttt{.runtime/} 文件夹以及 \texttt{assets/data/friends.db} 文件。这两个路径是本地数据库和身份信息的存放地。
\end{enumerate}

\subsection{Git 协作与 Commit 规范}
所有开发人员在提交代码时应遵循以下规范：
\begin{itemize}
    \item 目标主分支为 \texttt{main} 分支。
    \item 所有的 Commit 提交信息末尾，必须追加 Co-Authored 署名，格式如下：
    \begin{lstlisting}
    Fix socket connection issue in message relaying.
    
    Co-Authored-By: Claude <noreply@anthropic.com>
    \end{lstlisting}
\end{itemize}

\section{项目结构与模块职责}

\subsection{整体项目目录说明}
项目的核心文件和代码包按如下模块化组织：
\begin{lstlisting}
core/                           # 核心源代码根目录
  config/                       # 全局路径配置及跨平台路径适配
    paths.py                    # SQLite、缓存、文件接收目录的动态定位器
  frontend/                     # 前端界面逻辑
    views/                      # 核心视图页面
      discover.py               # 附近发现与好友请求页面
      friends.py                # 好友列表与管理页面
      chat.py                   # 即时聊天与分块文件传输页面
      moments.py                # P2P 朋友圈/北洋空间动态分享墙
      profile.py                # 个人信息及自动保存设置页面
    app.py                      # BeiyangApp 控制器，前端逻辑与后端服务的纽带
    theme.py                    # 统一的 Material 3 配色方案与设计系统
  backend/                      # 后端核心服务
    services/                   # 独立运行的网络与数据层服务
      social_runtime.py         # 系统生命周期总控、回调挂载与协调中心
      friend_db.py              # SQLite 数据库核心（包含好友、聊天记录、中继消息）
      udp_service.py            # 基于 UDP 的局域网主动扫描与被动发现服务
      connection_manager.py     # 基于 TCP Socket 连结池的底层双向通信与心跳管理
      message_service.py        # 消息中继机制、好友同步与通信路由解析
      file_transfer_state.py    # 传输任务状态监控器（断点状态、速度、进度）
      file_store.py             # 文件接收安全写控制与哈希完整性校验
      update_service.py         # 模拟软件版本更新中继发布服务
    shared/                     # 共享常量与网络协议包
      protocol.py               # 定义二进制协议头、握手和帧格式
      file_message.py           # 前后端共享传输状态数据帧结构
assets/                         # UI 静态资源（系统字体、内置头像、ICO 图标）
docs/                           # 系统相关文档目录
operations/                     # 运维与仿真测试工具包
  run_local_pair.ps1            # 双实例快速联调运行脚本
  run_random_users.ps1          # 大规模随机节点稳定性压力测试脚本
build-apk-clean.ps1             # 容器化 Android APK 零缓存编译打包脚本
docker-build-apk.sh             # Linux 下的 Docker 编译 Shell 脚本
pyproject.toml                  # Flet 项目打包依赖和版本配置文件
requirements.txt                # 独立部署依赖库清单
\end{lstlisting}

\subsection{系统分层设计原则}
本项目遵循严格的分层隔离设计：
\begin{enumerate}
    \item \textbf{前端视图层}：不直接操作网络 Socket，不直接连接 SQLite 数据库，所有数据变更均通过 \texttt{BeiyangApp} 提供的业务 API 执行。
    \item \textbf{后端服务层}：运行于独立线程中，提供高效的底层锁机制以确保数据安全性。各服务状态的变化通过事件回调（Callbacks）传递给应用管理器，进而驱动前端 UI 更新。
    \item \textbf{共享与配置层}：提供静态协议常量和安全路径定位，屏蔽操作系统（Windows / macOS / Linux / Android）的文件系统 and 底层接口差异。
\end{enumerate}

\section{功能特性与操作指南}

\subsection{附近的用户与好友请求机制}
\begin{itemize}
    \item \textbf{自动扫描}：进入“发现”页面后，应用会通过 UDP 自动扫描并展示当前局域网内的其他活动节点。
    \item \textbf{好友条件审核}：在“个人资料”中，您可以配置兴趣标签并设置“最低匹配标签数”。如果对方发送申请时的兴趣标签与您的匹配数少于此数值，系统会自动拒绝申请。
    \item \textbf{申请处理}：可在发现页的“好友请求”列表中点击“同意”或“拒绝”。支持在设置中开启“自动同意好友申请”开关。
\end{itemize}

\subsection{即时通信（Chat）与历史记录}
\begin{itemize}
    \item \textbf{单聊与群聊}：双击好友可发起点对点 TCP 聊天。系统同时支持创建局域网多点群聊，群聊元数据会在成员之间自动同步。
    \item \textbf{未读红点与滚动}：当有新消息且当前不处于该会话时，好友头像上会显示未读消息数红点。进入聊天后，窗口会自动平滑滚动到底部。
\end{itemize}

\subsection{文件传输、校验与断点续传}
\begin{itemize}
    \item \textbf{拖拽与卡片发送}：直接将文件拖入聊天窗口或点击“发送文件”按钮发送。
    \item \textbf{传输控制与断点续传}：大文件传输过程中支持随时“取消”。若传输中断，双击“重试”会自动获取本地已写入的字节位置，发起断点续传。
    \item \textbf{完成操作}：文件传输完成后，卡片上会自动提供“打开文件”（调用系统默认应用）、“定位目录”（在资源管理器中选中并高亮该文件）、“复制路径”按钮。如果是 ZIP 压缩包，还会提供“解压 ZIP”按钮。
\end{itemize}

\subsection{北洋空间动态分享墙（Moments）}
\begin{itemize}
    \item \textbf{动态发布}：进入“北洋空间”页面，您可以输入文字，并点击图片按钮添加一张本地图片作为附件，点击“发布动态”发布。
    \item \textbf{P2P 动态同步}：空间动态不依赖服务器，发布后将暂时保存在本地数据库中。当您与任何好友建立 TCP 连接时，系统会自动在后台将您以及您收集到的最新动态同步给对端。
    \item \textbf{动态互动}：您可以为好友的动态点赞或发表评论，这些互动信息同样会通过 P2P 链条同步给所有共同好友。
\end{itemize}

\section{核心技术设计与架构实现}

\subsection{基于 UDP 的设备发现协议}
用户启动应用后，UDP 服务（监听端口 \texttt{8890}）将自动启动并广播包含以下载荷的 JSON 帧：
\begin{lstlisting}
{
  "name": "用户昵称",
  "port": 7779,
  "action": "ping"
}
\end{lstlisting}
\textbf{发现机制流程}：
\begin{itemize}
    \item \textbf{周期性广播}：每隔 5 秒向局域网广播子网段发送一次 \texttt{ping} 帧。
    \item \textbf{主动扫描响应}：当收到新用户的 \texttt{ping} 帧时，本端会主动回应一个单播 \texttt{pong} 帧。双方获得对方的 IP 地址和 TCP 端口号，将对方列入“附近的用户”。
\end{itemize}

\subsection{基于 TCP 的点对点通信与协议帧设计}
两个已经加为好友的节点在进行通信前，将通过 TCP Socket 进行三次握手式状态验证，协议传输通过自定义二级帧结构解析：
\begin{enumerate}
    \item \textbf{版本握手}：建立连接后立即交换彼此的 \texttt{pyproject.toml} 中的软件版本 and UUID。
    \item \textbf{心跳维持（Heartbeat）}：连接保持期间，每隔 15 秒发送一次心跳空包以检测对端死活。若连续 3 次无回应则断开连接。
\end{enumerate}

\subsection{断点续传与安全文件传输设计}
文件分块传输（每块大小默认 \texttt{256KB}）的协议状态流程如图所示：
\begin{itemize}
    \item \textbf{握手阶段}：发送方提供文件全名、字节长度 and SHA-256 哈希值。接收方检索数据库，若发现有未传输完毕的旧任务记录，则返回当前已写入文件指针位置（例如 \texttt{Offset: 4194304} 字节），自动激活断点续传。
    \item \textbf{传输阶段}：发送方从指定 Offset 开始读取文件块并发送，接收方以追加模式写入目标路径文件。
    \item \textbf{校验阶段}：传输结束，接收方独立重新计算整机文件的 SHA-256 并同原始哈希进行匹配。若匹配成功，状态置为已完成，否则丢弃文件报错。
\end{itemize}

\subsection{离线消息与多跳网络中继协议}
对于暂时离线（UDP 未在线且 TCP 断开）的好友，消息发送将遵循有限跳数的八卦广播协议（Gossip-based Store-and-Forward Routing）：
\begin{itemize}
    \item 本地将该消息标记为“离线”状态并保存在本地数据库表 \texttt{pending\_messages} 中。
    \item 当检测到局域网中存在任意与本端和目标端同为好友的“在线第三方”时，本端会尝试将此消息包裹为 \texttt{RELAY\_MESSAGE} 并加密传输给中继节点。
    \item 中继节点作为“消息中继器”将此消息持久化到本地库中，最大中继跳数控制为 3 跳，在中继节点偶遇目标节点上线时最终将消息路由送达。
\end{itemize}

\subsection{SQLite 数据库架构设计}
所有本地的好友、关系、设置和消息记录均保存在本地 SQLite 数据库中。数据库包含以下核心表结构：
\begin{itemize}
    \item \texttt{friends}: 保存好友的 UUID、昵称、最近在线 IP、端口、头像名、个人简介及状态。
    \item \texttt{chat\_history}: 聊天消息记录表，包含消息唯一标识 \texttt{msg\_id}、发送方、接收方、消息文本、时间戳和状态。
    \item \texttt{pending\_messages}: 离线缓存消息表。
    \item \texttt{received\_msg\_ids}: 中继消息去重表，用于过滤已成功接收的中继消息，防止洪泛风暴。
    \item \texttt{groups} \& \texttt{group\_chat\_history}: 本地群聊及群聊历史表。
    \item \texttt{moments} \& \texttt{moment\_comments}: 北洋空间动态及互动评论表。
\end{itemize}

\subsection{线程安全与数据库排他锁}
由于界面渲染线程、TCP 接收线程、UDP 发现线程都会并发操作 SQLite，为了避免 SQLite 发生锁死错误，\texttt{friend\_db.py} 实现了严格的线程排他锁机制。所有写操作必须加锁：
\begin{lstlisting}[language=python]
def save_chat_message(self, from_name, to_name, content, timestamp, msg_id):
    with self._lock:
        cursor = self.conn.cursor()
        cursor.execute("INSERT INTO chat_history ...")
        self.conn.commit()
\end{lstlisting}

\subsection{自动保存草稿机制（Auto-save Draft Pattern）}
为了避免移动端和桌面端用户频繁点击保存按钮，个人资料与设置页的输入框采用了 Auto-save 机制。
在输入过程中，\texttt{on\_change} 回调会持续将输入值同步到内存中的草稿变量 \texttt{\_draft\_*}。当输入框失去焦点（\texttt{on\_blur}）或用户点击输入框外部（\texttt{on\_tap\_outside}）时，触发自动写入数据库。
\textbf{设计考量}：由于部分平台上 \texttt{on\_blur} 可能会先于最后一次 \texttt{on\_change} 完成触发，脚本在保存时优先读取内存草稿 \texttt{\_draft\_*}，其次回退读取控件值，彻底规避了时序竞态问题。

\section{桌面端安装与快速部署}

\subsection{基础 Python 环境部署}
推荐使用 Python 3.10.x - 3.12.x 环境运行：
\begin{lstlisting}[language=bash]
# 创建并激活环境
python -m venv .venv
.\.venv\Scripts\Activate.ps1
# 安装项目依赖
pip install -r requirements.txt
\end{lstlisting}

\subsection{Windows 双实例与仿真用户压力测试}
项目提供了完善的本端快速联调仿真脚本：
\begin{enumerate}
    \item \textbf{快速双实例联调}：在 PowerShell 下运行 \texttt{operations/run\_local\_pair.ps1}，该脚本会自动拉起 \texttt{Alice} 和 \texttt{Bob} 两个配置隔离的窗口进行收发测试。
    \item \textbf{大规模网络仿真}：运行 \texttt{run\_random\_users.ps1}：
    \begin{lstlisting}[language=bash]
    # 模拟启动 10 个后台隐藏节点（端口自动递增分配）以测试局域网发现和中继消息路由压力
    powershell -File operations/run_random_users.ps1 -Count 10 -NoWindow $true
    \end{lstlisting}
    测试完毕后可使用以下命令一键清理：
    \begin{lstlisting}[language=bash]
    Get-Process python | Stop-Process
    \end{lstlisting}
\end{enumerate}

\section{Android 平台编译与移动端适配}

\subsection{适配设计与兼容性考量}
Android 作为受限移动系统，相比桌面端在网络与文件系统有极大的限制。本项目在 \texttt{core} 中预设了针对 Android 的一系列底层机制：
\begin{table}[htbp]
  \centering
  \caption{移动端 Android 兼容性与解决方案表}
  \begin{tabular}{lll}
    \toprule
    限制分类 & Android 底层表现 & 解决方案 \\
    \midrule
    网络通信 & 锁屏/后台时 UDP 广播无法收发 & 使用 Jnius 调取系统 MulticastLock 服务 \\
    文件系统 & assets 目录只读，无法修改配置文件 & 自动将可写数据迁移到 app 专属私有存储目录 \\
    UI 接口 & 移动端无法导入 tkinter 模块 & 捕获异常，使用原生网页 FilePicker 替代 \\
    交互限制 & 无桌面级的 \texttt{os.startfile} 功能 & 通过 \texttt{am start} 调用 Android Intent 广播 \\
    资源文件 & 无法解析 ICO 格式图标 & 增加平台判断，加载 PNG 适应系统图标 \\
    \bottomrule
  \end{tabular}
\end{table}

\subsection{容器化无缓存 APK 编译命令操作}
打包脚本 \href{file:///E:/Alotofink/Python/Challenge/challenge3/build-apk-clean.ps1}{\texttt{build-apk-clean.ps1}} 已经被重构，彻底移除了可能导致旧代码复用的 Gradle 全局缓存机制。

要生成无缓存、最新版本的 APK，请执行以下命令：
\begin{lstlisting}[language=bash]
# 运行全新的防缓存编译
.\build-apk-clean.ps1
\end{lstlisting}
编译输出将位于项目根目录的 \texttt{build/apk/meeting-in-beiyang.apk} 下。

\section{开发者指南与测试运行}

\subsection{Flet 0.85.3 框架接口规范}
为了防止编写出因 Flet 版本升级导致不兼容的代码，所有开发人员需严格遵守以下控件参数定义：
\begin{itemize}
    \item \textbf{TextField}: 使用 \texttt{helper="..."} 来展示输入框下方的辅助提示文本（禁止使用旧版的 \texttt{helper\_text}）。
    \item \texttt{Dropdown}: 下拉选择事件必须绑定在 \texttt{on\_select} 事件上（禁止绑定在 \texttt{on\_change} 上）。
    \item \textbf{Switch}: 状态切换事件保持绑定在 \texttt{on\_change} 上。
    \item \textbf{FilePicker}: 在 Flet 0.85+ 中，FilePicker 属于 Service，必须添加到 \texttt{page.services.append(picker)} 中，不可再作为 Overlay 控件直接加入 \texttt{page.overlay}。
\end{itemize}

\subsection{运行本地单元测试}
在提交新的代码或分支前，请利用本地 Pytest 框架对数据库和中继机制进行一致性测试：
\begin{lstlisting}[language=bash]
# 运行测试套件
python -m pytest core/tests -q
\end{lstlisting}

\end{document}
